\preface

This book is an expanded version of my lecture notes at the Institute for Advanced Study in Mathematics (IASM) at Zhejiang University. My initial goal was modest: to write a self-contained reference for the participants in the lectures. But I soon realized that many results circulating in the area had never been rigorously proved anywhere in the literature. Each attempt to resolve these loose ends produced new ones, and the notes grew steadily until they became the present book.

\section*{Three pluripotential theories.}
The aim of this book is to present, from a unified point of view, three different but interrelated theories that all deserve to be called a \emph{global pluripotential theory}:
\begin{enumerate}
    \item\label{itm:preface:L7:1} the pluripotential theory on compact Kähler manifolds;
    \item\label{itm:preface:L7:2} the pluripotential theory on the Berkovich analytification of a projective variety;
    \item\label{itm:preface:L7:3} the toric pluripotential theory on a projective toric variety.
\end{enumerate}
The three theories are bound together by a common circle of ideas. Many results were originally proved in one setting and have analogues, sometimes evident and sometimes far from evident, in the others.

\section*{The central question: classifying singularities.}
Let us begin with the analytic side. Fix a compact Kähler manifold $X$. The central objects are the \emph{quasi-plurisubharmonic functions} on $X$ --- functions which are locally the sum of a smooth function and a plurisubharmonic function. We are mostly interested in their \emph{singularities}: the locus where $\varphi=-\infty$ and the rate at which $\varphi$ tends to $-\infty$ near this locus.

Singularities are not pathological exceptions; they are forced upon us by geometry. In typical applications, $X$ is the compactification of the moduli space of some class of geometric objects, with a Zariski-open locus $U\subseteq X$ parameterizing smooth objects. The natural metric on the associated polarizing line bundle is then smooth only on $U$, and the Grauert--Remmert extension theorem (\cref{thm:GRexten2}) extends it to all of $X$ at the cost of introducing singularities along $X\setminus U$. Understanding such singularities is therefore not an aesthetic refinement but a prerequisite for any moduli-theoretic application.

The local theory of these singularities is intricate, and a complete classification is essentially out of reach: locally, (quasi-)plurisubharmonic functions exhibit extremely complicated behavior, and the natural local invariants (Lelong numbers and multiplier ideal sheaves) capture only their rough features. For example, a function with log-log singularities has the same local invariants as a bounded one.

The picture changes drastically when one passes to the global setting. On a compact manifold, there are three natural equivalence relations on $\QPSH(X)$, of decreasing fineness:
\begin{enumerate}
    \item\label{itm:preface:L22:1} the \emph{singularity-type} relation $\sim$, identifying $\varphi,\psi$ whose difference $\varphi-\psi$ is bounded;
    \item\label{itm:preface:L22:2} the \emph{$P$-equivalence} $\sim_P$, two potentials with $\varphi\leq \psi$ are identified if the mixed non-pluripolar Monge--Amp\`ere masses
    \[
        \int_X \omega_\varphi^j\wedge \omega_\phi^{n-j} = \int_X \omega_\psi^j\wedge \omega_\phi^{n-j},\qquad j=0,\ldots,n,
    \]
    agree against every reference potential $\phi$. Here $\omega$ is a sufficiently large Kähler form. In other words, $\varphi$ and $\psi$ behave similarly with respect to the global non-pluripolar mass (\cref{def:Pmoresing});
    \item\label{itm:preface:L22:3} the \emph{$\mathcal{I}$-equivalence} $\sim_\mathcal{I}$, identifying $\varphi,\psi$ for which the multiplier ideal sheaves agree, $\mathcal{I}(k\varphi)=\mathcal{I}(k\psi)$ for every $k\in\mathbb{Z}_{>0}$ --- equivalently, $\varphi$ and $\psi$ have the same generic Lelong numbers along every prime divisor over every modification of $X$ (\cref{def:Imoresing}).
\end{enumerate}
The first relation records the finest information one can hope to capture by global invariants; $\sim_P$ retains the global \emph{mass} side, while $\sim_\mathcal{I}$ retains the algebraic-geometric side. They are studied systematically in \cref{chap:comp}.

A natural plan to classify singularities is the three-step decomposition:
\begin{enumerate}
    \item\label{itm:preface:L34:1} classify the $\mathcal{I}$-equivalence classes;
    \item\label{itm:preface:L34:2} classify the $P$-equivalence classes inside a given $\mathcal{I}$-class;
    \item\label{itm:preface:L34:3} classify the singularity types inside a given $P$-class.
\end{enumerate}
Step~3 has been the subject of much activity in the past decade under the name of \emph{pluripotential theory with prescribed singularities}, and is by now reasonably well understood. The first goal of this book is to give a complete answer to Step~1, and a satisfactory answer to Step~2 in the cases that arise in geometry.

\section*{$\mathcal{I}$-good singularities.}
Step~1 is the subject of \cref{chap:Igood}. We characterize $\mathcal{I}$-equivalence classes very explicitly through approximations by an explicit class of singularities, known as \emph{analytic singularities}; the analogous problem of approximating a plurisubharmonic function by ones with analytic singularities was studied in detail by Demailly.

The class of $\mathcal{I}$-good potentials is the central protagonist of this book. In a single equivalence theorem (\cref{thm:charIgoodasclosure}), $\mathcal{I}$-goodness admits three intuitive descriptions, each capturing a different facet of what it means for a singularity to be ``well-behaved.'' In the non-degenerate case, a quasi-plurisubharmonic function $\varphi$ with $\theta_{\varphi}=\theta+\ddc\varphi\geq 0$ for some closed smooth real $(1,1)$-form $\theta$ is $\mathcal{I}$-good if and only if any of the following hold:
\begin{enumerate}
    \item\label{itm:preface:L45:1} $\varphi$ is the limit, with respect to the $d_S$-pseudometric introduced in \cref{chap:comp}, of a sequence of potentials with analytic singularities. (\emph{Topological description}: $\mathcal{I}$-good potentials are precisely those reachable as limits of the explicit, well-understood class of analytic singularities.)
    \item\label{itm:preface:L45:2} The non-pluripolar Monge--Amp\`ere mass of $\varphi$ equals the algebraic volume introduced in \cref{def:volqpsh}:
    \[
        \int_X \theta_\varphi^n = \vol\theta_\varphi.
    \]
    (\emph{Mass--volume identity}: no mass is lost to pluripolar sets, so the analytic mass and the algebraic volume coincide.)
    \item\label{itm:preface:L45:3} The $P$-envelope and the $\mathcal{I}$-envelope of $\varphi$ introduced in \cref{chap:envelopes} agree:
    \[
        P_\theta[\varphi] = P_\theta[\varphi]_{\mathcal{I}}.
    \]
    (\emph{Analytic--algebraic envelope agreement}: the two natural least-singular representatives, one constructed from the $P$-equivalence relation and one from multiplier ideals, coincide.)
\end{enumerate}
The reader is invited to keep these three pictures in mind throughout the book; we will use whichever is most convenient at any given moment.

Step~2 is the more delicate one, and its difficulty turns out to be a feature, not a bug.

A given $\mathcal{I}$-equivalence class can consist of an enormous number of $P$-equivalence classes; constructing such examples is easy. By contrast, in the toric and non-Archimedean settings, Step~2 is essentially trivial: every $\mathcal{I}$-equivalence class consists of a single $P$-equivalence class. In the toric case, an $\mathcal{I}$- (or $P$-) equivalence class is a sub-convex body of the Newton body. In the non-Archimedean case, it is a homogeneous plurisubharmonic metric.

The discrepancy between the analytic and the toric/non-Archimedean theories suggests that, in the pluripotential theory on compact Kähler manifolds, a great many $\mathcal{I}$-classes are pathological, and that what one really wants is to single out within each $\mathcal{I}$-class a canonical $P$-class. The functions in this canonical class will be called \emph{$\mathcal{I}$-good}. The development of the theory of $\mathcal{I}$-good singularities, in \cref{chap:Igood}, is the technical core of this book. As the reader will see, almost every singularity that arises naturally in geometric applications is $\mathcal{I}$-good.

A useful guiding analogy is with measure theory. There exist many non-measurable functions, but unless one constructs one by hand using the axiom of choice, one only encounters measurable functions in practice. Similarly, although the space of non-$\mathcal{I}$-good singularities is non-empty, it is essentially invisible to geometric applications. Once one accepts the framework of $\mathcal{I}$-good singularities, Step~2 becomes essentially solved.

\section*{Tools, applications, conventions.}
The classification of singularities, in the form just described, is somewhat abstract. To put it to use, we develop two general techniques that allow induction on $\dim X$. For a prime divisor $Y$ in general position, the \emph{analytic Bertini theorem} relates quasi-plurisubharmonic functions on $X$ and on $Y$. For non-generic $Y$, the \emph{trace operator} associates with $\varphi$ a quasi-plurisubharmonic function on the normalization of $Y$, well-defined modulo $\mathcal{I}$-equivalence. These tools occupy \cref{chap:trace}.

The toric pluripotential theory on ample line bundles is mostly classical and is recalled, with full proofs, in \cref{chap:toric_ample}. The corresponding theory for big toric line bundles, which has never been written down in the literature, requires the partial Okounkov body machinery developed in \cref{chap:Okounkov} and is the content of \cref{chap:toric_big}.

The applications to non-Archimedean pluripotential theory occupy \cref{chap:NApotential}, building on the test-curve theory of \cref{chap:testcurve}. The convergence of partial Bergman kernels with arbitrary singular weight, extending the Berman--Boucksom--Witt Nyström theorem, is proved in \cref{chap:Bergman}.

\section*{Reading guide.}

This book is not an elementary introduction to pluripotential theory. It should not be the reader's first textbook in this field. The reader is expected to be comfortable with the basic pluripotential theory at the level of \cite{GZ17}, which is essentially the only prerequisite for the main body of the book.

This book is not ``a compilation of all the author's works''. Rather, it presents a new point of view on this theory, with new results and new insights. A reader who has gone through all my related papers is still encouraged to read at least part of this book.

This book is not an encyclopedia of pluripotential theory. For example, the reader would not find the definition of a pseudoconvex domain, nor the full details from the long research papers or books I depend on in the application part. These points irritated the referee, but I hope that these omissions are understandable, taking into account that this book is already over 500 pages long. Although I tried my best to make the book as self-contained as possible, certain proofs based on techniques far beyond the scope of this book are still omitted.

Results in this book are presented in the greatest possible generality. An expert reader may use this book as a convenient reference for the most general statements, and skip the proofs of the more general results if they are already familiar with them in a special case.

The reader is expected to read Part~II linearly, and consult the necessary preliminaries in Part~I when necessary.

A reader interested only in one application chapter may skip directly to it once the relevant Part~II prerequisites are in place. \cref{chap:toric_ample} is needed only for \cref{chap:toric_big}; it can otherwise be skipped without loss.

Some chapters in this book draw on more specialised background, recorded in the appendices: \cref{chap:convex} for convex analysis, \cref{chap:unibranch} for pluripotential theory on unibranch K\"ahler spaces, and \cref{chap:almostsg} for the almost-semigroup machinery underlying the partial Okounkov body construction.

In order to develop the theory of b-divisors as in \cref{chap:bdiv}, we need to take limits with respect to the net of all modifications of $X$. In general, this net does not admit a cofinal sequence.
As a consequence, nets are widely used in this book, and the reader is expected to be familiar with their basic properties. The books \cite{Wil70, Kel75} are good references for the basic theory.
It is essential to remember that the dominated convergence theorem and Fatou's lemma both fail for nets. A version of the monotone convergence theorem does hold, as in \cite[Proposition~7.12]{Fol99}.

I made a great effort to make sure that every single expression involving $(-\infty)-(-\infty)$ or $0\cdot \infty$ is explained. This point is often overlooked in the literature. In most cases, the fix is simple, but in some cases, it requires a non-trivial amount of work. Therefore a number of our proofs seem more complicated than in the original papers. Similarly, in the literature, many constants are only correct up to a constant involving $2$ or $\pi$. I tried to make sure that every constant in this book is correct, and hence there is an extensive discussion of the normalizations of various objects in \cref{chap:pshfunctions}, which again irritated the same referee. I hope that the reader will appreciate the effort and find the book more readable as a result.

Comments and questions are always welcome, and I will maintain an \href{https://mingchenxia.github.io/Lectures/Errata_SGPT.pdf}{errata page} on my homepage as long as I am still alive, no matter whether I remain in academia or not.

\vspace{\baselineskip}
\begin{flushright}\noindent
\hfill \textit{Mingchen Xia}\\
in Hangzhou, March 2024\\
\end{flushright}

\enlargethispage{14\baselineskip}
\vfill

\begin{large}
Contact information:

\textit{Mingchen Xia}, \textsc{\href{https://igp.ustc.edu.cn/main.htm}{Institute of Geometry and Physics}, USTC}\,\href{https://www.ustc.edu.cn/}{\raisebox{-0.25ex}{\includegraphics[height=1.2em]{images/ustc_emblem.png}}}\par\nopagebreak
\textit{Email address}: \texttt{xiamingchen2008@gmail.com}\par\nopagebreak
\textit{Homepage}: \href{https://mingchenxia.github.io/}{\raisebox{-0.2ex}{\includegraphics[height=1.1em]{images/rufa_icon.png}}}\,\url{https://mingchenxia.github.io/}.
\end{large}

\par\smallskip
\noindent\makebox[\textwidth][c]{%
    \includegraphics[width=0.75\textwidth]{images/rufa_xiaozhuan2.png}%
}
\par\vspace*{2.5\baselineskip}
